% BHU PYQ -- Transcribed & Corrected 100% Accurate
% Paper: DHAMD21: Tantra Vigyan (Indian Tantra Tradition)
% Exam: Under Graduate Semester II Examination 2024-25 (NEP)
% Category: Multidisciplinary Course (MDC)

\documentclass[11pt,a4paper]{article}
\usepackage[a4paper,top=2.5cm,bottom=2.5cm,left=2.5cm,right=2.5cm,headheight=14pt]{geometry}
\usepackage{amsmath,amssymb}
\usepackage[T1]{fontenc}
\usepackage[utf8]{inputenc}
\usepackage{lmodern,microtype,fancyhdr,enumitem,hyperref,lastpage,parskip}

\hypersetup{colorlinks=true,urlcolor=black,linkcolor=black,pdftitle={DHAMD21: Tantra Vigyan -- BHU Sem II 2024-25 (NEP MDC)}}
\pagestyle{fancy}\fancyhf{}
\fancyhead[L]{\small\textit{Banaras Hindu University}}
\fancyhead[R]{\small\textit{DHAMD21: Tantra Vigyan (MDC)}}
\fancyfoot[C]{\small Page \thepage\ of \pageref{LastPage}}

\newlist{parts}{enumerate}{2}
\setlist[parts,1]{label=(\alph*),leftmargin=*,itemsep=4pt,topsep=2pt}
\newcommand{\pts}[1]{\hfill\textit{[#1]}}

\begin{document}

\begin{center}
  {\large\bfseries BANARAS HINDU UNIVERSITY}\\[2pt]
  {\normalsize Under Graduate Semester II Examination 2024-25 (NEP)}\\[6pt]
  \rule{0.8\linewidth}{0.5pt}\\[6pt]
  {\large\bfseries MULTIDISCIPLINARY COURSE (MDC)}\\[4pt]
  {\large\bfseries Paper Code: DHAMD21 --- Tantra Vigyan (Indian Tantra Tradition)}\\[4pt]
  {\normalsize Time: 2:30 Hours\hfill Maximum Marks: 70}
\end{center}

\rule{\linewidth}{0.4pt}\smallskip
\noindent\textit{Instruction: Write your Roll No. at the top immediately on the receipt of this question paper.}\\[2pt]
\noindent\textit{Note: Answer any FIVE questions. All questions carry equal marks (14 marks each).}
\bigskip

\noindent\textbf{Question 1}\pts{14} \\
Provide a detailed introduction to the Indian Tantra Tradition (\textit{Bhartiya Tantra Parampara}).

\medskip\rule{\linewidth}{0.2pt}\medskip

\noindent\textbf{Question 2}\pts{14} \\
Highlight the significance and fundamental principles of Tantra Shastra in Indian philosophy.

\medskip\rule{\linewidth}{0.2pt}\medskip

\noindent\textbf{Question 3}\pts{14} \\
Describe the concept and practices of Shaktism and Tantra Sadhana.

\medskip\rule{\linewidth}{0.2pt}\medskip

\noindent\textbf{Question 4}\pts{14} \\
Discuss the interrelationship between Tantra and Modern Science (\textit{Tantra \& Vigyan}).

\medskip\rule{\linewidth}{0.2pt}\medskip

\noindent\textbf{Question 5}\pts{14} \\
Write a concise essay on Mantra Science (\textit{Mantra Vigyan}) and its psychological/spiritual impact.

\medskip\rule{\linewidth}{0.2pt}\medskip

\noindent\textbf{Question 6}\pts{14} \\
Elaborate on the significance of Yantra and Mandala Shastra in sacred geometry and Tantric rituals.

\medskip\rule{\linewidth}{0.2pt}\medskip

\noindent\textbf{Question 7}\pts{14} \\
Describe the nature and philosophical symbolism of the Shiva and Shakti principles (\textit{Shiva-Shakti Tattva}).

\vfill\rule{\linewidth}{0.4pt}
\begin{center}\small
\href{https://bhu-exam.vercel.app/}{Click here to visit our website}\\[4pt]
\href{https://me-aryan.vercel.app/}{Click here to visit our contributor}
\end{center}

\end{document}
